\section{Представящ слой и потребителско взаимодействие}

Back-end логиката на една система за електронна търговия има стойност само ако бъде експонирана чрез последователен и използваем потребителски интерфейс. В разглеждания проект визуалният слой е реализиран чрез Thymeleaf шаблони, допълнени от CSS, Bootstrap базирани компоненти и ограничено количество JavaScript. Този избор е в пълно съответствие с общата архитектурна постановка на приложението: интерфейсът е сървърно рендиран, което улеснява интеграцията между модели, валидация и HTML форми.

Настоящата глава разглежда организацията на шаблоните, ролята на навигацията, локализацията, формите и основните UX решения. Анализът се базира върху реалните HTML шаблони в директорията \texttt{src/main/resources/templates} и статичните ресурси в \texttt{src/main/resources/static}.

\subsection{Организация на шаблоните}

Визуалният слой е структуриран около набор от отделни Thymeleaf шаблони, всеки от които отговаря на конкретен потребителски екран. Сред тях се открояват следните страници:

\begin{itemize}[leftmargin=*]
    \item \texttt{index.html} -- начална страница и входна точка за неавтентикирани потребители;
    \item \texttt{login.html} и \texttt{register.html} -- форми за удостоверяване и регистрация;
    \item \texttt{profile.html} -- екран за преглед и редакция на профила;
    \item \texttt{add-company.html} -- форма за добавяне на фирма;
    \item \texttt{add-product.html} -- форма за добавяне на продукт;
    \item \texttt{products-catalog.html} -- основен каталожен екран;
    \item \texttt{details.html} -- детайлен изглед на конкретен продукт;
    \item \texttt{shopping-cart.html} -- количка и форма за checkout;
    \item \texttt{history.html} -- история на направените поръчки;
    \item \texttt{users.html}, \texttt{companies.html}, \texttt{delivery.html} и error шаблоните -- допълнителни визуални компоненти на системата.
\end{itemize}

Този набор от шаблони е достатъчен, за да обхване целия потребителски цикъл в приложението. Важно е да се отбележи, че визуалният слой не е структуриран чрез централен layout механизъм или фрагменти, а по-скоро чрез отделни самостоятелни HTML файлове. Това прави проекта лесен за локално проследяване, но води и до значително дублиране на елементи като навигационната лента.

\begin{table}[H]
    \centering
    \caption{Съответствие между шаблони и основни потребителски функции}
    \label{tab:template-mapping}
    \begin{tabularx}{\textwidth}{L{3.3cm}YY}
        \toprule
        Шаблон & Основна цел & Потребителски контекст \\
        \midrule
        \texttt{index.html} & Представяне на начална страница и достъп до основни действия & Гост или неавтентикиран потребител \\
        \texttt{register.html} & Създаване на нов акаунт & Гост \\
        \texttt{login.html} & Вход в системата & Гост или непотвърден потребител \\
        \texttt{profile.html} & Преглед и редакция на данни, списък с фирми & Удостоверен потребител \\
        \texttt{add-company.html} & Въвеждане на фирма & Удостоверен потребител \\
        \texttt{add-product.html} & Въвеждане на продукт & Удостоверен потребител \\
        \texttt{products-catalog.html} & Каталожен преглед и филтриране & Удостоверен потребител \\
        \texttt{details.html} & Избор на конкретен продукт и количество & Удостоверен потребител \\
        \texttt{shopping-cart.html} & Обобщение на количка и checkout & Удостоверен потребител \\
        \texttt{history.html} & Преглед и изчистване на историята & Удостоверен потребител \\
        \bottomrule
    \end{tabularx}
\end{table}

Таблица \ref{tab:template-mapping} показва, че интерфейсът е построен около ясни и тематично отделени екрани. Това е добра практика, защото всяка страница има разпознаваема функция и лесно може да бъде обвързана с конкретен контролерен маршрут.

\subsection{Навигация и общи интерфейсни елементи}

Един от най-видимите повторяеми елементи в шаблоните е навигационната лента. Тя присъства в почти всички екрани и включва връзки към регистрация, вход, профил, каталог, добавяне на продукт, история, количка, изход и смяна на езика. От гледна точка на потребителската логика това е силно решение: основните действия са достъпни от почти всеки екран и навигационната структура остава позната.

В същото време архитектурно решението има недостатък. Понеже навигацията се повтаря ръчно в множество шаблони, всяка бъдеща промяна трябва да бъде приложена на много места. При по-голямо приложение това би увеличило вероятността от несъответствия и би направило поддръжката по-трудна. По-зрял подход би бил използването на Thymeleaf fragments или общ layout шаблон. Независимо от това, за настоящия проект текущият вариант е лесен за проследяване и не затруднява разбирането на страниците.

Допълнителен общ елемент е визуалният фон, който в много страници използва фиксирано изображение и сходна цветова схема. Това придава визуална последователност, макар и реализирана чрез inline стилове, а не чрез централен CSS дизайн слой.

\subsection{Thymeleaf като механизъм за свързване на данни и интерфейс}

Thymeleaf \cite{thymeleaf} играе централна роля в представящия слой. Чрез атрибути като \texttt{th:text}, \texttt{th:field}, \texttt{th:each}, \texttt{th:if}, \texttt{th:href} и \texttt{th:action} шаблоните могат директно да използват данните, подготвени от контролерите. Това значително улеснява връзката между back-end слоя и HTML формите.

Особено полезно е използването на \texttt{th:field} при формите. По този начин се постига естествено свързване между binding модел и HTML полета, както и възможност за визуализиране на състояние при грешка. Например при формите за добавяне на фирма, продукт или редакция на профил, полетата се обвързват директно с модела, а условната визуализация на грешки използва \texttt{\#fields.hasErrors(...)}.

\begin{lstlisting}[caption={Типичен Thymeleaf фрагмент за форма},label={lst:thymeleaf-form}]
<form th:action="@{/products/add-product}"
      th:object="${productAddBindingModel}"
      th:method="POST">
    <label for="name">Name</label>
    <input th:field="*{name}" id="name" type="text" class="form-control">
    <p th:if="${#fields.hasErrors('name')}" class="errors alert alert-danger">
        Name must be between 5 and 20 characters.
    </p>
</form>
\end{lstlisting}

Листинг \ref{lst:thymeleaf-form} показва как се реализира типичен формов сценарий. От една страна, HTML остава достатъчно четим. От друга страна, връзката с Java модела е непосредствена. Именно в това се състои едно от основните предимства на сървърно рендирания подход при Spring MVC приложения.

\subsection{Форми и обработка на потребителски вход}

Почти всички съществени действия в системата са реализирани чрез форми. Това включва регистрация, вход, редакция на профил, добавяне на фирма, добавяне на продукт, добавяне на количество в количката и checkout процеса. Формите са реализирани в стил, характерен за Bootstrap \cite{bootstrap}, с класове като \texttt{form-group}, \texttt{form-control}, \texttt{btn} и други визуални помощници.

В профилния екран се вижда сравнително богата композиция: форма за лични данни, визуализиране на текущите стойности и таблица с компаниите на потребителя. Това е добър пример за страница, която съчетава редакция и преглед на свързана информация. Формата за добавяне на фирма пък е по-компактна и логически разделя текстовите от числовите полета. Формата за добавяне на продукт е най-добре структурирана по отношение на домейнните данни, тъй като включва избор от enum категории и избор на фирма от предварително зареден списък.

От UX гледна точка позитивен аспект е, че при невалиден вход потребителят получава визуална обратна връзка чрез alert елементи. Това подобрява използваемостта и прави формовия поток по-ясен. Ограничение е, че съобщенията не са напълно стандартизирани на всички места, а част от тях са твърдо кодирани в самите HTML файлове.

\subsection{Каталог и продуктови детайли}

Каталогът е един от най-характерните екрани в системата. Той е реализиран като комбинация от навигационна зона, странично меню с категории, ценови филтър и галерия от продуктови карти. Визуалният модел се опира на Bootstrap, но включва и значително количество inline CSS, описващ оформлението на елементите, размерите на картите, поведението на менюто и отзивчивостта на страницата.

Съществено UX качество на каталога е, че всеки продукт се визуализира с изображение, име, категория и цена, а бутонът за детайлен преглед отвежда потребителя към следващата стъпка от взаимодействието. В допълнение клиентският филтър по цена придава усещане за по-динамичен интерфейс, без да се изисква пълноценна SPA архитектура.

Страницата с детайли на продукта е по-концентрирана върху единичен обект. Тя представя изображение, име, цена и наличност, а формата за въвеждане на количество позволява директно добавяне към количката. От гледна точка на потребителския поток това е правилно решение: каталогът служи за избор, а страницата за детайли -- за вземане на решение и заявяване на покупка.

\subsection{Количка и checkout екран}

Екранът за количка комбинира две взаимосвързани цели: визуализация на текущото съдържание и подготовка на финализирането на поръчката. Визуалната композиция разделя съдържанието на две части: списък с избраните продукти и summary панел със суми, доставка, адрес и бутон за checkout.

Това решение е потребителски удачно, тъй като показва едновременно детайла на избраните артикули и обобщената финансова рамка на поръчката. Формата за адрес е вградена директно в summary панела, което намалява броя на преходните екрани. От друга страна, самата страница е относително тежка като inline CSS и в production контекст би било по-удачно стиловете да се централизирани във външен ресурс.

\subsection{Профил и история на поръчките}

Профилният екран е сред най-богатите откъм комбинирано съдържание. Той събира както form-based редакция на основните потребителски полета, така и индиректно насърчава управлението на фирмите чрез страничен панел със списък от съществуващи компании и връзка за добавяне на нова. Това е добра UX практика, тъй като поставя логически свързаните операции в един и същ контекст.

Историята на поръчките е реализирана с тъмен табличен интерфейс, който се различава визуално от останалите екрани. Това прави страницата разпознаваема и внушава архивен характер на информацията. Таблицата показва идентификатор на поръчката, имейл, адрес, дата и цена. Допълнително е наличен бутон за изчистване на историята. От потребителска гледна точка този екран изпълнява коректно ролята си, макар че в production система биха били желани филтриране, търсене и по-богат контекст за отделните поръчки.

\subsection{Локализация и двуезичен интерфейс}

Поддръжката на два езика е реализирана чрез \texttt{messages.properties} и \texttt{messages\_bg.properties}, заедно с \texttt{LocaleChangeInterceptor}. В навигационната лента са добавени преки връзки \texttt{?lang=en} и \texttt{?lang=bg}, които позволяват смяна на езика чрез параметър в заявката. Това е прост, но ефективен модел за локализация в сървърно рендирана система.

От анализа на шаблоните се вижда, че голяма част от основните надписи използват message ключове чрез \texttt{\#\{...\}}. Това е добро решение, защото държи езиковите ресурси извън самите шаблони. Съществуват обаче и частично локализирани или напълно твърдо кодирани елементи. Например част от текстовете за категории, някои помощни изрази и някои съобщения за грешка не са напълно синхронизирани между двата езика. Това означава, че локализацията е реално налична, но все още не е напълно завършена.

\subsection{Използване на външни библиотеки и клиентски ресурси}

Визуалният слой използва Bootstrap, Font Awesome, jQuery и jQuery UI \cite{bootstrap,jquery}. Това е видно както от шаблоните за каталог и детайли, така и от количеството CDN препратки в HTML файловете. От гледна точка на бързото разработване това е разбираем избор: готовите UI библиотеки позволяват по-бързо постигане на визуално приемлив резултат.

Наред с положителните страни тук има и няколко наблюдения. Първо, в различни шаблони се използват различни версии на Bootstrap. Второ, част от ресурсите са включени едновременно от локални файлове и от CDN. Трето, голямо количество стилове се дефинира inline във всеки отделен шаблон, вместо да се поддържа по-строга централизирана front-end организация. Това не е критичен проблем за малък проект, но е важен индикатор за бъдещи възможности за подобрение.

\subsection{UX силни страни}

Въпреки опростения си характер, интерфейсът има няколко съществени предимства:

\begin{itemize}[leftmargin=*]
    \item основните действия са видими и достъпни от навигацията;
    \item формите са ясно отделени по функции и показват грешки при входни проблеми;
    \item каталогът е визуално разпознаваем и поддържа филтриране по категория и цена;
    \item checkout екранът обединява продуктова информация и обобщение на поръчката;
    \item профилът и историята предлагат логично групирани потребителски сценарии.
\end{itemize}

Тези характеристики правят системата използваема и подходяща за демонстрация пред научен ръководител, комисия или крайни потребители в контекста на дипломната разработка.

\subsection{Ограничения на представящия слой}

Както и при back-end частта, представящият слой има своите ограничения:

\begin{itemize}[leftmargin=*]
    \item дублиране на навигационна логика и HTML блокове между шаблоните;
    \item смесване на външни CDN ресурси, локални файлове и inline CSS;
    \item частична, но не пълна локализация на интерфейса;
    \item отсъствие на централен layout или fragment система;
    \item наличие на твърдо кодирани визуални и текстови елементи, които намаляват гъвкавостта.
\end{itemize}

Тези ограничения обаче не отменят факта, че интерфейсът е функционален и достатъчно изразителен за целите на проекта. По-скоро те очертават естествена траектория за последващо професионализиране на front-end слоя.

\subsection{Обобщение}

Представящият слой на проекта демонстрира последователен сървърно рендиран подход, базиран на Thymeleaf, Bootstrap и локализирани ресурси. Шаблоните покриват пълния жизнен цикъл на потребителските действия в системата и са добре обвързани с back-end логиката. От инженерна гледна точка основните им предимства са простотата, директната интеграция с Spring MVC и достатъчно ясната визуална организация.

Същевременно анализът разкрива и типични ограничения на малък учебен проект: дублиране, непълна централизираност на стила и частична локализация. Тези наблюдения са важни, защото показват, че интерфейсът не е само визуален завършек, а реален слой със собствени архитектурни качества и слабости. Следващата глава разглежда структурата на проекта, конфигурацията на средата и генерираните build артефакти, които осигуряват operational аспект на системата.
